\subsection{Invariant Subspaces, Reducing Pairs, and Quotient Operators}

Throughout this subsection, $\mathcal{V}$ is a vector space over $\mathbb{F}$.\\
In statements involving inner products, adjoints, or unitary matrices, $\mathbb{F}\in\{\mathbb{R},\mathbb{C}\}$.

\begin{definition} \label{def:invariant-subspace} \hlt{Invariant Subspace}\\
Let $T\in\LL(\mathcal{V})$. A subspace $\mathcal{U}\leq\mathcal{V}$ is \hlt{invariant under $T$} iff $T\mathcal{U}\subseteq\mathcal{U}$. The set of all $T$-invariant subspaces is $\operatorname{Inv}(T)\defeq\{\mathcal{U}\leq\mathcal{V}:T\mathcal{U}\subseteq\mathcal{U}\}$.
\end{definition}

\begin{definition} \label{def:restriction-and-quotient-operators} \hlt{Restriction and Quotient Operators}\\
Let $T\in\LL(\mathcal{V})$ and $\mathcal{U}\in\operatorname{Inv}(T)$.
\begin{enumerate}[label=(\roman*)]
\setlength{\itemsep}{0pt}
\item The \hlt{restriction operator} is $T|_{\mathcal{U}}:\mathcal{U}\to\mathcal{U}$, $T|_{\mathcal{U}}(\mathbf{u})\defeq T\mathbf{u}$.
\item The \hlt{quotient operator} is $T/\mathcal{U}:\mathcal{V}/\mathcal{U}\to\mathcal{V}/\mathcal{U}$, $(T/\mathcal{U})(\mathbf{v}+\mathcal{U})\defeq T\mathbf{v}+\mathcal{U}$.
\end{enumerate}
\end{definition}

\begin{proposition} \label{prop:restriction-and-quotient-operators-well-defined} \hlt{Restriction and Quotient Operators Are Well-Defined}\\
Let $T\in\LL(\mathcal{V})$ and $\mathcal{U}\in\operatorname{Inv}(T)$.
\begin{enumerate}[label=(\roman*)]
\setlength{\itemsep}{0pt}
\item $T|_{\mathcal{U}}\in\LL(\mathcal{U})$.
\item $T/\mathcal{U}\in\LL(\mathcal{V}/\mathcal{U})$.
\end{enumerate}
\begin{proof}
\begin{enumerate}[label=(\roman*)]
\setlength{\itemsep}{0pt}
\item The inclusion $T\mathcal{U}\subseteq\mathcal{U}$ holds by Definition \ref{def:invariant-subspace}, so $T|_{\mathcal{U}}$ has codomain $\mathcal{U}$.\\
For $a,b\in\mathbb{F}$ and $\mathbf{u},\mathbf{v}\in\mathcal{U}$, $T|_{\mathcal{U}}(a\mathbf{u}+b\mathbf{v})=aT|_{\mathcal{U}}\mathbf{u}+bT|_{\mathcal{U}}\mathbf{v}$ by the linearity of $T$.\\
Therefore $T|_{\mathcal{U}}\in\LL(\mathcal{U})$ by the codomain inclusion and linearity identity.
\item Suppose $\mathbf{v}+\mathcal{U}=\mathbf{w}+\mathcal{U}$. Then $\mathbf{v}-\mathbf{w}\in\mathcal{U}$ by Proposition \ref{prop:parallel-affine-subsets-equal-or-disjoint}.\\
The relation $T\mathbf{v}-T\mathbf{w}=T(\mathbf{v}-\mathbf{w})\in\mathcal{U}$ holds by linearity and invariance.\\
Consequently, $T\mathbf{v}+\mathcal{U}=T\mathbf{w}+\mathcal{U}$ by Proposition \ref{prop:parallel-affine-subsets-equal-or-disjoint}, so the quotient rule is independent of the representative.\\
Set $S\defeq T/\mathcal{U}$ and $[\mathbf{x}]\defeq\mathbf{x}+\mathcal{U}$. For $a,b\in\mathbb{F}$ and $\mathbf{v},\mathbf{w}\in\mathcal{V}$, $S(a[\mathbf{v}]+b[\mathbf{w}])=[T(a\mathbf{v}+b\mathbf{w})]$ by Proposition \ref{prop:quotient-space-structure}(ii).\\
The identity $[T(a\mathbf{v}+b\mathbf{w})]=[aT\mathbf{v}+bT\mathbf{w}]=aS[\mathbf{v}]+bS[\mathbf{w}]$ holds by linearity of $T$ and the quotient rule.\\
Therefore $T/\mathcal{U}\in\LL(\mathcal{V}/\mathcal{U})$ by representative independence and the two linearity identities.
\end{enumerate}
\hlt{Structural Commentary}: Invariance closes the restriction codomain and transports equality between quotient representatives.\\
\hlt{Key Insight}: Preserving $\mathcal{U}$ also preserves differences modulo $\mathcal{U}$.
\end{proof}
\end{proposition}

\begin{proposition} \label{prop:one-dimensional-invariant-subspaces-and-eigenvectors} \hlt{One-Dimensional Invariant Subspaces and Eigenvectors}\\
Let $T\in\LL(\mathcal{V})$ and $\mathbf{v}\in\mathcal{V}\setminus\{\mathbf{0}\}$. Then $\operatorname{span}(\mathbf{v})\in\operatorname{Inv}(T)$ iff $\mathbf{v}$ is an eigenvector of $T$.
\begin{proof}
\begin{enumerate}[label=\roman*.]
\setlength{\itemsep}{0pt}
\item[$(\Rightarrow)$] Suppose $\operatorname{span}(\mathbf{v})\in\operatorname{Inv}(T)$. The relation $(\exists\,\lambda\in\mathbb{F})\,T\mathbf{v}=\lambda\mathbf{v}$ holds by $T\mathbf{v}\in\operatorname{span}(\mathbf{v})$.\\
Therefore $\mathbf{v}$ is an eigenvector of $T$ by $\mathbf{v}\neq\mathbf{0}$ and Definition \ref{def:eigenvalue-eigenvector}.
\item[$(\Leftarrow)$] Suppose $\mathbf{v}$ is an eigenvector of $T$. Choose $\lambda\in\mathbb{F}$ satisfying $T\mathbf{v}=\lambda\mathbf{v}$ by Definition \ref{def:eigenvalue-eigenvector}.\\
The inclusion $T(\operatorname{span}(\mathbf{v}))\subseteq\operatorname{span}(\mathbf{v})$ holds by $(\forall\,a\in\mathbb{F})\,T(a\mathbf{v})=a\lambda\mathbf{v}$. Hence $\operatorname{span}(\mathbf{v})\in\operatorname{Inv}(T)$ by Definition \ref{def:invariant-subspace}.
\end{enumerate}
\hlt{Structural Commentary}: The two directions translate line invariance into the image relation for its generator.\\
\hlt{Key Insight}: A line is invariant exactly when its generator is rescaled.
\end{proof}
\end{proposition}

Figure \ref{fig:eigenline-scaling} illustrates Proposition \ref{prop:one-dimensional-invariant-subspaces-and-eigenvectors}: $T$ rescales $\mathbf{v}$ within $\operatorname{span}(\mathbf{v})$.\\
The proposition, rather than the diagram, owns the equivalence.

\begin{illustrationfigure}[H]
  \centering
  \begin{tikzpicture}[
  x=0.88cm,
  y=0.86cm,
  >=Stealth,
  every node/.style={font=\small},
  axis/.style={->, gray!65, line width=0.45pt},
  eigline/.style={dartmouthgreen, line width=1.15pt, dashed},
  vec/.style={->, black, line width=1.05pt},
  imagevec/.style={->, magenta, line width=1.15pt},
  ghost/.style={gray!40, line width=0.45pt}
]
  \coordinate (O) at (0,0);
  \coordinate (v) at (1.45,0.82);
  \coordinate (Tv) at (3.05,1.72);

  \draw[axis] (-3.1,0) -- (3.55,0);
  \draw[axis] (0,-1.55) -- (0,2.15);

  \draw[eigline] (-2.75,-1.55) -- (3.25,1.84);
  \node[dartmouthgreen, anchor=south east] at (-1.4,-0.8) {$E_{\lambda}(T)$};

  \draw[vec] ($(O)+(0,-0.08)$) -- ($(v)+(0,-0.08)$)
    node[midway, below right=1pt] {$\mathbf{v}$};
  \draw[imagevec] ($(O)+(0,0.08)$) -- ($(Tv)+(0,0.08)$)
    node[midway, above=2pt] {$T\mathbf{v}=\lambda\mathbf{v}$};

  \draw[ghost, ->] (v) .. controls (1.95,1.15) and (2.35,1.35) .. (Tv);
  \node[anchor=north west] at (0.15,-1.05) {$\mathbf{v}\neq\mathbf{0}$};
  \node[anchor=south west, jazzberryjam] at (2.2,1.9) {same line};
\end{tikzpicture}

  \caption{An eigenvector and its image lie on the same invariant line.}
  \label{fig:eigenline-scaling}
\end{illustrationfigure}

\begin{proposition} \label{prop:invariant-subspaces-and-block-form} \hlt{Invariant Subspaces and Block Form}\\
Let $\mathcal{V}$ be finite-dimensional, let $T\in\LL(\mathcal{V})$, suppose $\mathcal{V}=\mathcal{U}\oplus\mathcal{W}$, and let $\beta$ be a basis adapted to this decomposition.
\begin{enumerate}[label=(\roman*)]
\setlength{\itemsep}{0pt}
\item $\mathcal{U}\in\operatorname{Inv}(T)\Longrightarrow[T]_{\beta\leftarrow\beta}$ is block upper triangular.
\item $\mathcal{U},\mathcal{W}\in\operatorname{Inv}(T)\Longrightarrow[T]_{\beta\leftarrow\beta}$ is block diagonal.
\end{enumerate}
\begin{proof}
Write $\beta=(\mathbf{u}_1,\ldots,\mathbf{u}_p,\mathbf{w}_1,\ldots,\mathbf{w}_q)$, where $(\mathbf{u}_1,\ldots,\mathbf{u}_p)$ and $(\mathbf{w}_1,\ldots,\mathbf{w}_q)$ are bases of $\mathcal{U}$ and $\mathcal{W}$, respectively.
\begin{enumerate}[label=(\roman*)]
\setlength{\itemsep}{0pt}
\item Suppose $\mathcal{U}\in\operatorname{Inv}(T)$. The relation $(\forall\,j\in\{1,\ldots,p\})\,T\mathbf{u}_j\in\mathcal{U}$ holds by invariance.\\
Therefore the coordinates along $\mathbf{w}_1,\ldots,\mathbf{w}_q$ vanish, forming the zero lower-left block.
\item Suppose also that $\mathcal{W}\in\operatorname{Inv}(T)$. The relation $(\forall\,j\in\{1,\ldots,q\})\,T\mathbf{w}_j\in\mathcal{W}$ holds by invariance.\\
Therefore the coordinates along $\mathbf{u}_1,\ldots,\mathbf{u}_p$ vanish, forming the zero upper-right block.
\end{enumerate}
\hlt{Structural Commentary}: Adapted-basis columns convert invariance of each summand into a vanishing off-diagonal block.\\
\hlt{Key Insight}: Each invariant summand eliminates the block pointing out of it.
\end{proof}
\end{proposition}

\begin{definition} \label{def:reducing-pair-and-direct-sum-operator} \hlt{Reducing Pair}\\
Let $T\in\LL(\mathcal{V})$. A pair $(\mathcal{U},\mathcal{W})$ of subspaces \hlt{reduces $T$} iff $\mathcal{V}=\mathcal{U}\oplus\mathcal{W}$ and $\mathcal{U},\mathcal{W}\in\operatorname{Inv}(T)$.\\
For such a pair, $T=T|_{\mathcal{U}}\oplus T|_{\mathcal{W}}$ means $(\forall\,(\mathbf{u},\mathbf{w})\in\mathcal{U}\times\mathcal{W})\,T(\mathbf{u}+\mathbf{w})=T|_{\mathcal{U}}\mathbf{u}+T|_{\mathcal{W}}\mathbf{w}$.
\end{definition}

\begin{theorem} \label{thm:adjoint-invariant-subspace-orthogonal-complement} \hlt{Adjoints and Invariant Orthogonal Complements}\\
Let $T\in\LL(\mathcal{V})$ on a finite-dimensional inner-product space and let $\mathcal{U}\leq\mathcal{V}$.
\begin{enumerate}[label=(\roman*)]
\setlength{\itemsep}{0pt}
\item $\mathcal{U}\in\operatorname{Inv}(T)\Longleftrightarrow\mathcal{U}^{\perp}\in\operatorname{Inv}(T^*)$.
\item $(\mathcal{U},\mathcal{U}^{\perp})$ reduces $T\Longleftrightarrow\mathcal{U}\in\operatorname{Inv}(T)\cap\operatorname{Inv}(T^*)$.
\item If $(\mathcal{U},\mathcal{U}^{\perp})$ reduces $T$, then $(T|_{\mathcal{U}})^*=(T^*)|_{\mathcal{U}}$.
\end{enumerate}
\begin{proof}
\begin{enumerate}[label=(\roman*)]
\setlength{\itemsep}{0pt}
\item[$(i)$] ($\Rightarrow$) Suppose $\mathcal{U}\in\operatorname{Inv}(T)$. For $\mathbf{z}\in\mathcal{U}^{\perp}$ and $\mathbf{u}\in\mathcal{U}$, $\innerprod{T^*\mathbf{z},\mathbf{u}}=\innerprod{\mathbf{z},T\mathbf{u}}=0$ by conjugate symmetry, Definition \ref{def:adjoint-linear-map}, and $T\mathbf{u}\in\mathcal{U}$.\\
Thus $T^*\mathbf{z}\in\mathcal{U}^{\perp}$, and $\mathcal{U}^{\perp}\in\operatorname{Inv}(T^*)$ by Definition \ref{def:invariant-subspace}.
\item[$(i)$] ($\Leftarrow$) Suppose $\mathcal{U}^{\perp}\in\operatorname{Inv}(T^*)$. If $\mathbf{u}\in\mathcal{U}$ and $\mathbf{z}\in\mathcal{U}^{\perp}$, then $\innerprod{T\mathbf{u},\mathbf{z}}=\innerprod{\mathbf{u},T^*\mathbf{z}}=0$ by Definition \ref{def:adjoint-linear-map}.\\
Therefore $T\mathbf{u}\in(\mathcal{U}^{\perp})^{\perp}=\mathcal{U}$ by Corollary \ref{cor:double-orthogonal-complement}, so $\mathcal{U}\in\operatorname{Inv}(T)$ by Definition \ref{def:invariant-subspace}.
\item[$(ii)$] ($\Rightarrow$) Suppose $(\mathcal{U},\mathcal{U}^{\perp})$ reduces $T$. Then $\mathcal{U}\in\operatorname{Inv}(T)$ and $\mathcal{U}^{\perp}\in\operatorname{Inv}(T)$ by Definition \ref{def:reducing-pair-and-direct-sum-operator}.\\
Thus $\mathcal{U}\in\operatorname{Inv}(T^*)$ by (i) for $T^*$ and Proposition \ref{prop:basic-adjoint-identities}(iii).
\item[$(ii)$] ($\Leftarrow$) Suppose $\mathcal{U}\in\operatorname{Inv}(T)\cap\operatorname{Inv}(T^*)$. Then $\mathcal{U}^{\perp}\in\operatorname{Inv}(T)$ by (i) for $T^*$ and Proposition \ref{prop:basic-adjoint-identities}(iii).\\
The identity $\mathcal{V}=\mathcal{U}\oplus\mathcal{U}^{\perp}$ holds by Theorem \ref{thm:orthogonal-direct-sum-decomposition}. Hence $(\mathcal{U},\mathcal{U}^{\perp})$ reduces $T$ by Definition \ref{def:reducing-pair-and-direct-sum-operator}.
\item[$(iii)$] Suppose $(\mathcal{U},\mathcal{U}^{\perp})$ reduces $T$. Then $\mathcal{U}\in\operatorname{Inv}(T)\cap\operatorname{Inv}(T^*)$ by (ii), and $T|_{\mathcal{U}},(T^*)|_{\mathcal{U}}\in\LL(\mathcal{U})$ by Proposition \ref{prop:restriction-and-quotient-operators-well-defined}(i).\\
For $\mathbf{u},\mathbf{v}\in\mathcal{U}$, $\innerprod[\mathcal{U}]{T|_{\mathcal{U}}\mathbf{u},\mathbf{v}}=\innerprod[\mathcal{U}]{\mathbf{u},(T^*)|_{\mathcal{U}}\mathbf{v}}$ by Definition \ref{def:adjoint-linear-map}. Identity $(T|_{\mathcal{U}})^*=(T^*)|_{\mathcal{U}}$ follows by Definition \ref{def:adjoint-linear-map}, Lemma \ref{lem:adjoint-unique-if-exists}.
\end{enumerate}
\hlt{Structural Commentary}: Adjoint duality transfers invariance to $\mathcal{U}^{\perp}$; orthogonal decomposition then characterizes reduction.\\
\hlt{Key Insight}: Reducing $\mathcal{U}$ is joint invariance under $T$ and $T^*$.
\end{proof}
\end{theorem}

\begin{proposition} \label{prop:reducing-pairs-and-algebraic-projections} \hlt{Reducing Pairs and Projections}\\
Let $T\in\LL(\mathcal{V})$, let $\mathcal{V}=\mathcal{U}\oplus\mathcal{W}$, and set $\rho\defeq\rho_{\mathcal{U},\mathcal{W}}$. Then $(\mathcal{U},\mathcal{W})$ reduces $T$ iff $T\rho=\rho T$.
\begin{proof}
\begin{enumerate}[label=\roman*.]
\setlength{\itemsep}{0pt}
\item[$(\Rightarrow)$] Suppose $(\mathcal{U},\mathcal{W})$ reduces $T$. For $\mathbf{v}=\mathbf{u}+\mathbf{w}$ with $(\mathbf{u},\mathbf{w})\in\mathcal{U}\times\mathcal{W}$, $(T\rho)\mathbf{v}=T\mathbf{u}=(\rho T)\mathbf{v}$ by $T\mathbf{u}\in\mathcal{U}$ and $T\mathbf{w}\in\mathcal{W}$.\\
Therefore $T\rho=\rho T$ by pointwise equality.
\item[$(\Leftarrow)$] Suppose $T\rho=\rho T$. For $\mathbf{u}\in\mathcal{U}$, $\rho T\mathbf{u}=T\rho\mathbf{u}=T\mathbf{u}$, so $T\mathbf{u}\in\operatorname{range}(\rho)=\mathcal{U}$ by Proposition \ref{prop:algebraic-projection-range-kernel-idempotent}(iii).\\
For $\mathbf{w}\in\mathcal{W}$, $\rho T\mathbf{w}=T\rho\mathbf{w}=\mathbf{0}$, so $T\mathbf{w}\in\ker(\rho)=\mathcal{W}$ by Proposition \ref{prop:algebraic-projection-range-kernel-idempotent}(iii). Hence $(\mathcal{U},\mathcal{W})$ reduces $T$ by Definition \ref{def:reducing-pair-and-direct-sum-operator}.
\end{enumerate}
\hlt{Structural Commentary}: Pointwise commutation preserves the range and kernel summands of the algebraic projection.\\
\hlt{Key Insight}: Commutation with $\rho$ is simultaneous invariance of its range and kernel.
\end{proof}
\end{proposition}

\begin{proposition} \label{prop:invariant-subspace-projection-criterion} \hlt{Invariant Subspaces and Algebraic Projections}\\
Let $T\in\LL(\mathcal{V})$, let $\mathcal{V}=\mathcal{U}\oplus\mathcal{W}$, and set $\rho\defeq\rho_{\mathcal{U},\mathcal{W}}$.
\begin{enumerate}[label=(\roman*)]
\setlength{\itemsep}{0pt}
\item $\mathcal{U}\in\operatorname{Inv}(T)\Longleftrightarrow\rho T\rho=T\rho$.
\item If $\mathcal{U}\in\operatorname{Inv}(T)$, then $(\forall\,\mathcal{Z}\leq\mathcal{V})\,[\mathcal{V}=\mathcal{U}\oplus\mathcal{Z}\Longrightarrow\rho_{\mathcal{U},\mathcal{Z}}T\rho_{\mathcal{U},\mathcal{Z}}=T\rho_{\mathcal{U},\mathcal{Z}}]$.
\end{enumerate}
\begin{proof}
\begin{enumerate}[label=(\roman*)]
\setlength{\itemsep}{0pt}
\item[$(i)$] ($\Rightarrow$) Suppose $\mathcal{U}\in\operatorname{Inv}(T)$, and write $\mathbf{v}=\mathbf{u}+\mathbf{w}$ with $(\mathbf{u},\mathbf{w})\in\mathcal{U}\times\mathcal{W}$.\\
The identity $(\rho T\rho)\mathbf{v}=\rho T\mathbf{u}=T\mathbf{u}=(T\rho)\mathbf{v}$ holds by $T\mathbf{u}\in\mathcal{U}$ and Definition \ref{def:projection-onto-subspace-along-complement}. Therefore $\rho T\rho=T\rho$ by pointwise equality.
\item[$(i)$] ($\Leftarrow$) Suppose $\rho T\rho=T\rho$. If $\mathbf{u}\in\mathcal{U}$, then $\rho T\mathbf{u}=\rho T\rho\mathbf{u}=T\rho\mathbf{u}=T\mathbf{u}$.\\
Hence $T\mathbf{u}\in\operatorname{range}(\rho)=\mathcal{U}$ by Proposition \ref{prop:algebraic-projection-range-kernel-idempotent}(iii), so $\mathcal{U}\in\operatorname{Inv}(T)$ by Definition \ref{def:invariant-subspace}.
\item[$(ii)$] Suppose $\mathcal{U}\in\operatorname{Inv}(T)$ and fix $\mathcal{Z}\leq\mathcal{V}$ satisfying $\mathcal{V}=\mathcal{U}\oplus\mathcal{Z}$. Set $\sigma\defeq\rho_{\mathcal{U},\mathcal{Z}}$.\\
The identity $\sigma T\sigma=T\sigma$ holds by the forward implication in (i) for $\mathcal{V}=\mathcal{U}\oplus\mathcal{Z}$. Claim (ii) follows by the arbitrary choice of $\mathcal{Z}$.
\end{enumerate}
\hlt{Structural Commentary}: The compression identity tests range preservation and remains independent of the chosen complement.\\
\hlt{Key Insight}: $\rho T\rho=T\rho$ forbids $T$ from leaving $\operatorname{range}(\rho)$.
\end{proof}
\end{proposition}

\begin{theorem} \label{thm:matrix-invariant-subspace-projection-block-criteria} \hlt{Invariant Subspaces, Blocks, and Projections}\\
Let $n\in\mathbb{Z}_{>0}$, $A\in\operatorname{Mat}(n,n,\mathbb{F})$, and $\mathcal{U}\leq\mathbb{F}^n$, where $\dim(\mathcal{U})=r$ and $1\leq r\leq n-1$.\\
Let $Q=[Q_1\ Q_2]$ be unitary, where the columns of $Q_1$ and $Q_2$ are orthonormal bases of $\mathcal{U}$ and $\mathcal{U}^{\perp}$, respectively.\\
Set $P\defeq Q_1Q_1^*$. The following are equivalent:
\begin{enumerate}[label=(\roman*)]
\setlength{\itemsep}{0pt}
\item $A\mathcal{U}\subseteq\mathcal{U}$.
\item The lower-left $(n-r)\times r$ block of $Q^*AQ$ is zero.
\item $PAP=AP$.
\end{enumerate}
\begin{proof}
The identities $Q_1^*Q_1=I_r$, $Q_2^*Q_1=0$, $P^*=P$, and $P^2=Q_1(Q_1^*Q_1)Q_1^*=P$ hold by unitarity of $Q$.\\
For $\mathbf{x}\in\mathbb{F}^n$ and $\mathbf{c}\in\mathbb{F}^r$, $P\mathbf{x}\in\operatorname{col}(Q_1)$ and $P(Q_1\mathbf{c})=Q_1\mathbf{c}$, so $\operatorname{range}(P)=\operatorname{col}(Q_1)=\mathcal{U}$.\\
For $\mathbf{z}\in\mathcal{U}^{\perp}$, $Q_1^*\mathbf{z}=\mathbf{0}$ and $P\mathbf{z}=\mathbf{0}$. Therefore $P=\rho_{\mathcal{U},\mathcal{U}^{\perp}}$ by Theorem \ref{thm:orthogonal-direct-sum-decomposition} and Definition \ref{def:projection-onto-subspace-along-complement}.\\
The matrix $P$ is an orthogonal projection matrix by Definition \ref{def:orthogonal-projection-matrix}.\\
Set $M\defeq Q^*AQ=\left[\begin{smallmatrix}B&X\\Y&C\end{smallmatrix}\right]$, so $Y=Q_2^*AQ_1$, and set $E\defeq Q^*PQ=I_r\oplus0_{n-r}$.
\begin{enumerate}[label=(\roman*)]
\setlength{\itemsep}{0pt}
\item[$(i)$] $\Rightarrow (ii)$ Suppose $A\mathcal{U}\subseteq\mathcal{U}$. Every column of $AQ_1$ belongs to $\mathcal{U}=(\mathcal{U}^{\perp})^{\perp}$ by Corollary \ref{cor:double-orthogonal-complement}.\\
Since the columns of $Q_2$ span $\mathcal{U}^{\perp}$, $Y=Q_2^*AQ_1=0$.
\item[$(ii)$] $\Rightarrow (iii)$ Suppose $Y=0$. Then $EME=\left[\begin{smallmatrix}B&0\\0&0\end{smallmatrix}\right]=ME$, so $Q^*(PAP)Q=Q^*(AP)Q$.\\
Therefore $PAP=AP$ by multiplication by $Q$ and $Q^*$.
\item[$(iii)$] $\Rightarrow (i)$ Suppose $PAP=AP$. Since $P$ is the algebraic projection onto $\mathcal{U}$ along $\mathcal{U}^{\perp}$, $A\mathcal{U}\subseteq\mathcal{U}$ by Proposition \ref{prop:invariant-subspace-projection-criterion}(i).
\end{enumerate}
\hlt{Structural Commentary}: Unitary coordinates identify the escape block; projection compression tests the same range preservation.\\
\hlt{Key Insight}: The lower-left block records precisely the component leaving $\mathcal{U}$.
\end{proof}
\end{theorem}

\begin{corollary} \label{cor:reducing-subspace-projection-commutation} \hlt{Reducing Subspaces and Commuting Projections}\\
Let $n\in\mathbb{Z}_{>0}$, $A\in\operatorname{Mat}(n,n,\mathbb{F})$, and $\mathcal{U}\leq\mathbb{F}^n$, where $\dim(\mathcal{U})=r$ and $1\leq r\leq n-1$.\\
Let $Q=[Q_1\ Q_2]$ be unitary with $\operatorname{col}(Q_1)=\mathcal{U}$ and $\operatorname{col}(Q_2)=\mathcal{U}^{\perp}$. Set $P\defeq Q_1Q_1^*$. The following are equivalent:
\begin{enumerate}[label=(\roman*)]
\setlength{\itemsep}{0pt}
\item $A\mathcal{U}\subseteq\mathcal{U}$ and $A^*\mathcal{U}\subseteq\mathcal{U}$.
\item $Q^*AQ$ is block diagonal relative to $\mathbb{F}^r\oplus\mathbb{F}^{n-r}$.
\item $PA=AP$.
\end{enumerate}
\begin{proof}
Set $M\defeq Q^*AQ=\left[\begin{smallmatrix}B&X\\Y&C\end{smallmatrix}\right]$ and $E\defeq Q^*PQ=I_r\oplus0_{n-r}$. Then $Q^*A^*Q=M^*$, $Q^*(PA)Q=EM$, and $Q^*(AP)Q=ME$.
\begin{enumerate}[label=(\roman*)]
\setlength{\itemsep}{0pt}
\item[$(i)$] $\Rightarrow (ii)$ Suppose $A\mathcal{U}\subseteq\mathcal{U}$ and $A^*\mathcal{U}\subseteq\mathcal{U}$. Then $Y=0$ by Theorem \ref{thm:matrix-invariant-subspace-projection-block-criteria}, (i)$\Rightarrow$(ii), for $A$.\\
The relation $X^*=0$ holds by the same implication for $A^*$ because $X^*$ is the lower-left block of $M^*$.\\
Therefore $X=Y=0$, so $M=B\oplus C$.
\item[$(ii)$] $\Rightarrow (iii)$ Suppose $M=B\oplus C$. Then $EM=ME$ by block multiplication, and $PA=AP$ by conjugation with $Q$.
\item[$(iii)$] $\Rightarrow (i)$ Suppose $PA=AP$. Then $EM=ME$. The identities $EM=\left[\begin{smallmatrix}B&X\\0&0\end{smallmatrix}\right]$ and $ME=\left[\begin{smallmatrix}B&0\\Y&0\end{smallmatrix}\right]$ imply $X=Y=0$.\\
Therefore $A\mathcal{U}\subseteq\mathcal{U}$ by Theorem \ref{thm:matrix-invariant-subspace-projection-block-criteria}, (ii)$\Rightarrow$(i), for $A$.\\
Since $X^*=0$ is the lower-left block of $M^*$, $A^*\mathcal{U}\subseteq\mathcal{U}$ by the same implication for $A^*$.
\end{enumerate}
\hlt{Structural Commentary}: The invariant-block criterion for $A$ and $A^*$ eliminates opposite off-diagonal blocks.\\
\hlt{Key Insight}: Projection commutation is equivalent to simultaneous $A$- and $A^*$-invariance.
\end{proof}
\end{corollary}

Figure \ref{fig:invariant-projection-block-form} summarizes the preceding results: invariance removes the lower-left block.\\
Projection commutation removes both off-diagonal blocks.\\
Proposition \ref{prop:invariant-subspaces-and-block-form}, Theorem \ref{thm:matrix-invariant-subspace-projection-block-criteria}, and Corollary \ref{cor:reducing-subspace-projection-commutation} own these claims.

\begin{illustrationfigure}[H]
  \centering
  \begin{tikzpicture}[
  >=Stealth,
  every node/.style={font=\small},
  space/.style={draw=blue!60!black, fill=blue!7, rounded corners=3pt, minimum width=2.05cm, minimum height=0.78cm, align=center},
  mat/.style={draw=gray!60, fill=gray!5, rounded corners=2pt, minimum width=2.05cm, minimum height=1.35cm, align=center},
  arrow/.style={->, line width=0.85pt, blue!70!black},
  weak/.style={->, line width=0.75pt, gray!65}
]
  \node[space] (U) at (0,0.72) {$\mathcal{U}$};
  \node[space, fill=teal!7, draw=teal!60!black] (Perp) at (0,-0.72) {$\mathcal{U}^{\perp}$};

  \node[mat] (Inv) at (3.75,0.72) {$\begin{bmatrix}B&X\\0&C\end{bmatrix}$};
  \node[mat] (Red) at (3.75,-0.72) {$\begin{bmatrix}B&0\\0&C\end{bmatrix}$};

  \draw[arrow] (U.east) -- (Inv.west);
  \draw[weak] (Perp.east) -- (Inv.west);
  \draw[arrow] (Inv.south) -- (Red.north);
\end{tikzpicture}

  \caption{Adapted coordinates represent invariance by upper block form and reduction by block diagonal form.}
  \label{fig:invariant-projection-block-form}
\end{illustrationfigure}

\begin{perspective} \label{per:restriction-quotient-perspective} \hlt{Two Smaller Views of an Operator}\\
Let $T\in\LL(\mathcal{V})$ and $\mathcal{U}\in\operatorname{Inv}(T)$. Definition \ref{def:restriction-and-quotient-operators} separates the action inside $\mathcal{U}$ from the action modulo $\mathcal{U}$.\\
Proposition \ref{prop:restriction-and-quotient-operators-well-defined} validates both operators. In an adapted block matrix from Proposition \ref{prop:invariant-subspaces-and-block-form}(i), they determine the two diagonal blocks but not the upper-right block. Thus restriction and quotient simplify the study of $T$ without necessarily determining $T$.
\end{perspective}
